\section{本章小结}

本章作为数理统计的基础，介绍了几个基本概念和3种常见的分布。

~

{\bf 总体}和{\bf 样本}

总体是所有考察对象的集合，这种“所有”往往带有不可穷尽，所以实际中我们通过获取部分对象，即样本$\left( X_1,X_2,\cdots ,X_n \right) $，考察总体的性质。
样本是一个$n$维随机变量。

{\bf 统计量}和{\bf 抽样分布}

统计量$g\left( X_1,X_2,\cdots ,X_n \right) $是样本的函数，一个多元随机变量函数，其值是必须能通过样本完全确定，才能用来估计总体。
其次，统计量也是一个随机变量，有自己的分布，称为抽样分布，具体形态由总体、统计量的函数形式$g$、样本的容量$n$共同决定。
好的抽样分布能反映总体的情况，常见的有卡方分布、t分布、F分布。




